%To make simple comments in the text
\newcommand{\com}[1]{\textcolor{red}{#1}}
%To define the vector command (which will just make math FAT)
\renewcommand{\vec}[1]{\mathbf{#1}}
%To make any chemical formulas look like chemical formulas
\newcommand*\chem[1]{\ensuremath{\mathrm{#1}}}
\DeclareSIUnit\Molar{\textsc{M}}
\DeclareSIUnit\molar{\textsc{M}}
% To make some chemicals properly hyphenated
% \hyphenation{
%     er-go-no-mic,% example
%     tri-me-thoxy-silyl,
%     me-tha-cry-late,
%     lu-ti-dine,
% }

% I really don't care about hyphenation in the bibliography, and some names are just not hyphenated, super annoying. So force the matter here.
\hyphenation{S-c-i-o-r-t-i-n-o}
\hyphenation{S-c-h-w-a-r-t-z}

% command for etal and roman numerals
\newcommand{\etal}{\textit{et al.~}} % usage \etal{}

\newcommand{\croman}[1]
{\MakeUppercase{\romannumeral #1}}

\newenvironment{captioncont}[1]{
  \begin{tcolorbox}[colback=gray!5, colframe=gray!30, boxrule=0.2pt, sharp corners, left=6pt, right=6pt, top=4pt, bottom=4pt]
  \small
  Figure~\ref{#1} \textit{caption (continued).}
}{\end{tcolorbox}}
\newcommand{\tred}[1]{\textcolor{red}{(#1)}}
\newcommand{\suppfig}[0]{Supplementary Fig.~}
\newcommand{\supptable}[0]{Supplementary Tab.~}



\newcommand{\chapterbar}{
  \pgfmathparse{\thechapter}
  \let\chapnum\pgfmathresult
  \ifnum\chapnum>0
    \ifnum\chapnum<7
      \checkoddpage
      \ifoddpage
        \pgfmathparse{\colourarray[\chapnum]}
        \edef\chapclr{\pgfmathresult}
        \begin{tikzpicture}[remember picture,overlay]
          \pgfmathsetmacro{\yoffset}{6 - ((\chapnum - 1) * (2.5))}
          \node[
            anchor=center,
            inner sep=2pt,
            font=\sffamily\bfseries\color{white},
            text width=1.5cm,
            align=center,
            fill=\chapclr,
            rounded corners=3pt,
            drop shadow,
            minimum height=1cm
          ] at ([xshift=-0.4cm, yshift=\yoffset cm]current page.east)
          {\chapnum};
        \end{tikzpicture}
      \fi
    \fi
  \fi
}



% for chapter bar to be a circle


% \newcommand{\chapterbar}{
%   \ifnum\value{chapter}>0
%     \checkoddpage
%     \ifoddpage
%       \pgfmathparse{\colourarray[\thechapter]}
%       \edef\chapclr{\pgfmathresult}
%       \begin{tikzpicture}[remember picture,overlay]
%         \pgfmathsetmacro{\yoffset}{6 - ((\thechapter - 1) * (14 / 7))}
%         \node[
%           circle,
%           anchor=center,
%           inner sep=2pt,
%           minimum size=1.4cm,
%           font=\sffamily\bfseries\color{white},
%           fill=\chapclr,
%           drop shadow,
%           align=center
%         ] at ([xshift=0.2cm, yshift=\yoffset cm]current page.east)
%         {\thechapter};
%       \end{tikzpicture}
%     \fi
%   \fi
% }



\AddEverypageHook{\chapterbar}

%chapter text color for figures alphabets. figure 2c-d, c-d will have the same color as the chapter
\newcommand{\chaptc}[1]{%
  \textcolor{\chapclr}{#1}%
}
% remove the header only on the first page of the chapter
\DeclareNewPageStyleByLayers{noheadrule}{
  scrheadings.foot.scrlayer % keep footer
}
%\deffootnote{1.5em}{1em}{\makebox[1.5em][l]{\thefootnotemark}}
\deffootnote{0em}{0em}{\thefootnotemark\ }
\setkomafont{footnote}{\footnotesize}
%\renewcommand{\footnoterule}{\kern-3pt\hrule width 0.4\columnwidth\kern 2.6pt}
\renewcommand{\footnoterule}{\vspace{2cm}\kern-3pt\hrule width 0.4\columnwidth\kern 2.6pt}
\setcapindent{2pt}  % remove the indent in the caption coming from the hanging indent in koma script
\DeclareSIUnit{\angstrom}{\textup{\AA}}
